\section*{Chapitre \ref{ch:s1:earnings} --- Résultats}

\begin{solution}{exo:s1:earnings:1}
$1.5 \times 1.2~\% = 1.8~\%$ sur soixante jours~; par unité de surprise $1.2~\%/60 = 2.0$ points de base par jour.
\end{solution}

\begin{solution}{exo:s1:earnings:2}
Le poids échangé vaut deux fois la rotation à sens unique~: $2 \times 65.2 \times 0.001 = 13.0~\%$ par an pour le livre à cinq jours, $2 \times 6.6 \times 0.001 = 1.3~\%$ pour
le livre à soixante jours. Le rendement brut du livre à cinq jours, 8.3~\% par an, ne couvre pas ses coûts.
\end{solution}

\begin{solution}{exo:s1:earnings:3}
Environ $0.05 \times 6.94\% = 0.35$ points, proche de 0.33, que montre l'\omterm{def:s1:earnings:event}{étude d'événement} (le mouvement contient aussi des composantes de marché et de secteur, qui ne sont pas inversées).
\end{solution}

\begin{solution}{exo:s1:earnings:4}
Le rendement du jour de l'annonce est le saut lié à la surprise plus le bruit propre à ce jour, de sorte qu'il mesure la surprise avec erreur~; la surprise publiée est la
cause injectée elle-même. Sur données réelles, la surprise publiée se mesure par rapport à un consensus qui peut être périmé ou biaisé, et le rendement reflète tout ce qui
a été annoncé (prévisions de la direction, marges), si bien que le rendement peut porter une information absente du chiffre principal. Combiner les deux est courant.
\end{solution}

\begin{solution}{exo:s1:earnings:5}
La date est fixée et publiée à l'avance~; la surprise ne devient publique qu'à l'annonce. Un backtest doit savoir si chaque annonce est tombée
avant l'ouverture, pendant la séance ou après la clôture, et n'agir qu'à partir du premier cours qui suit~: une annonce faite après la clôture ne peut pas être négociée au
cours de clôture de ce jour.
\end{solution}

\begin{solution}{exo:s1:earnings:6}
Sur cinq jours de cotation, dont un jour d'annonce de variance environ neuf fois celle d'un jour ordinaire ($3.02^2$), la variance vaut $4 + 9 = 13$ journées ordinaires
contre 5 sans l'événement~: la volatilité implicite de l'option couvrant l'annonce vaut environ $\sqrt{13/5} = 1.61$ fois celle d'une option qui ne
la couvre pas, et retombe une fois l'événement passé.
\end{solution}

\begin{solution}{exo:s1:earnings:7}
Les ratios de Sharpe valent 2.18 avant et 0.75 après la rupture~; 4.8~\% par an avant et 1.4~\% après. Une baisse de cette ampleur dans une série de volatilité d'environ 2~\% par
an demande plusieurs années pour se distinguer du bruit avec les tests de rupture (livre~7, chapitre~13)~; attendre la certitude revient à gérer pendant des années une stratégie affaiblie.
\end{solution}

\begin{solution}{exo:s1:earnings:8}
La dérive est concentrée dans les actions peu liquides~: Chordia et ses coauteurs trouvent 0.04~\% par mois dans les actions les plus liquides contre 2.43~\% dans les moins liquides, et
Martineau n'en trouve aucune dans les grandes actions depuis 2006. Un backtest sur l'ensemble des actions dit peu de chose du S\&P 500~; il faut le lancer sur cet univers, avec les coûts, et sur des
années récentes.
\end{solution}

\begin{solution}{pb:s1:earnings:1}
\begin{enumerate}
\item Rendements anormaux moyens en \omterm{def:s1:earnings:event}{temps d'événement}, sur une \omterm{def:s1:earnings:event}{fenêtre d'événement} autour des annonces~; la dérive est la poursuite des rendements anormaux dans le sens de la surprise.
\item Les surprises passées et les rendements passés prédisent chacun de fortes dérives une fois neutralisé l'effet de l'autre~; les prévisions des analystes réagissent elles aussi lentement.
\item Une surprise normale centrée réduite, un saut de trois volatilités spécifiques quotidiennes par unité et 1.2~\% de dérive par unité sur soixante jours~; quatre annonces par an
(1.59~\% des jours-actions).
\item Au-dessus de $+1$ écart-type~: 6.94~\% le jour même, 7.75~\% à soixante jours~; en dessous de $-1$~: $-6.93~\%$ et $-8.56~\%$.
\item Positions longues sur les titres dont la surprise dépasse un écart-type, positions courtes sur ceux dont elle est inférieure à moins un écart-type, avec un délai après l'annonce pour une fenêtre de détention,
exposition brute égale à un.
\item À la clôture de l'annonce~: $-3.72$, 0.02, 2.03~; le lendemain~: $-0.77$, 1.51, 2.66.
\item Le lendemain, une partie du saut s'inverse (0.33 points pour le groupe positif).
\item Chaque position paie un aller-retour complet pour quelques points de base de dérive par jour~: des coûts de 13~\% par an à cinq jours.
\item La surprise (2.66 contre 2.03 à soixante jours)~: le rendement mêle la surprise et le bruit.
\item Rendements plus élevés autour des annonces programmées, 9.9~\% par an annualisés chez Savor et Wilson~; aucun sur le marché synthétique (1.8 point de base par jour, $t =
0.59$).
\item 3.02 fois les mouvements ordinaires~; les options couvrant la date valorisent le saut, si bien que leur volatilité implicite est plus élevée.
\item Acheter les révisions à la hausse des prévisions, et trader sur les prévisions de la direction et les avertissements sur résultats~; historiques des prévisions telles qu'elles étaient connues à chaque date, et communiqués de presse horodatés.
\item \textbf{Résultat nommé.} Ratios de Sharpe après coûts de $-0.77$, 1.51 et 2.66 pour des durées de détention de 5, 20 et 60 jours (achat le lendemain
de l'annonce)~; après que la dérive a perdu 70~\% de son ampleur, le ratio de Sharpe du livre à soixante jours sur les années 7 à 10 vaut 0.75 contre 3.24.
\item La dérive est faible, voire absente, dans les actions liquides (0.04~\% par mois) et absente des grandes actions depuis 2006.
\item Dans les petites actions peu liquides, où les coûts la rongent, et dans les signaux que le marché digère plus lentement (prévisions de la direction, révisions, texte).
\item Suivre la dérive réalisée par unité de surprise, sa trajectoire en \omterm{def:s1:earnings:event}{étude d'événement} et le rendement du livre face à un test de rupture.
\item Horodatages avant ou après séance, consensus périmé, résultats retraités et dates d'annonce déplacées après coup.
\item Le livre de dérive à soixante jours sur les titres liquides où elle survit, de taille modeste~; ou la \omterm{def:s1:earnings:premium}{prime d'annonce}, traitée comme prime de risque.
\item Qu'un rendement concentré autour d'événements puisse être une rémunération du risque, et pas seulement une erreur.
\item Positions prises autour des annonces programmées, à parier sur le degré d'absorption de l'information par les cours.
\end{enumerate}
\end{solution}

\begin{solution}{iq:s1:earnings:1}
Rendements anormaux qui se poursuivent dans le sens d'une surprise de résultats pendant des semaines après l'annonce. Explications~: les investisseurs sous-réagissent et actualisent lentement~;
limites à l'arbitrage dans les actions peu liquides où il est le plus fort~; une prime de risque que la littérature n'a pas tranchée.
\end{solution}

\begin{solution}{iq:s1:earnings:2}
Résultats réels tels que publiés pour la première fois et consensus tel qu'il se présentait avant l'annonce, tous deux horodatés~; la surprise rapportée au cours ou à l'historique
des surprises passées (SUE). Les erreurs possibles~: les résultats retraités remplaçant les premières publications, les instantanés de consensus pris après l'annonce, les heures d'annonce
mal datées, et chiffres par action ajustés des divisions d'actions mêlés à des chiffres non ajustés.
\end{solution}

\begin{solution}{iq:s1:earnings:3}
Tester formellement la rupture, en estimer l'ampleur après la rupture, la mesurer là où le livre peut trader (titres liquides, coûts inclus), et se demander pourquoi elle s'est affaiblie. Si
le rendement après rupture dépasse encore les coûts avec une marge suffisante, conserver une allocation plus faible~; sinon, arrêter.
\end{solution}

\begin{solution}{iq:s1:earnings:4}
Comparer le mouvement implicite des options couvrant la date (à partir de l'écart entre leur variance implicite et celle des options voisines) avec une
prévision du mouvement effectif~; acheter des straddles couverts en delta lorsque la prévision dépasse le mouvement implicite de plus que les coûts, les vendre dans le cas inverse.
\end{solution}

\begin{solution}{iq:s1:earnings:5}
Risque de saut à chaque annonce, concentration sur les titres annonçant cette semaine-là, regroupement des annonces en pleine saison des résultats, et risque de données (horodatages, retraitements).
\end{solution}

\begin{solution}{iq:s1:earnings:6}
Avec des positions détenues $H$ jours puis remplacées, il y a $252/H$ allers-retours par an et par unité de capital. Pour $H \le 60$, le rendement net vaut
$(252/H)(dH/60 - 2c) = 252\,(d/60 - 2c/H)$, croissant en $H$~; pour $H > 60$, il vaut $(252/H)(d - 2c)$, décroissant. L'optimum consiste à détenir pendant toute la dérive,
$H = 60$.
\end{solution}
